\begin{centering}
\Large{\textbf{RESUMEN}}

\end{centering}

Los modelos de texto a voz basados en redes neuronales profundas dependen de manera crítica de la cantidad y la calidad de sus datos de entrenamiento, lo que representa una barrera significativa para idiomas y variedades dialectales de bajos recursos, como el español de Argentina. Frente a esta limitación, los conjuntos de datos recolectados de fuentes heterogéneas constituyen una alternativa atractiva, aunque requieren cadenas de preprocesamiento que filtren y estandaricen el material. En este trabajo se evaluó, mediante parámetros objetivos y subjetivos, el impacto de dichas cadenas de procesamiento sobre el conjuntos de datos y su efecto en el entrenamiento de modelos de TTS.

Para ello se desarrolló una cadena de preprocesamiento modular y reproducible, compuesta por etapas de detección de actividad de voz, realce de señal y filtrado por calidad perceptual no intrusiva, y se la aplicó bajo múltiples configuraciones, contrastando distintos algoritmos de denoising (DeepFilterNet y Demucs), modelos de estimación de calidad (NISQA y DNSMOS) y umbrales de filtrado. La evaluación se articuló en torno a una métrica compuesta que integra reducción del corpus, calidad de señal, condiciones acústicas y preservación de las características del hablante. De forma complementaria, se entrenó un modelo de estimación de densidad basado en difusión probabilística sobre un corpus profesional de referencia, y se validaron los hallazgos en tareas de clonación de voz mediante modelos de síntesis \textit{zero-shot}.

Los resultados evidenciaron una divergencia fundamental entre ambos criterios de evaluación: mientras las métricas objetivas intrusivas premiaron la supresión agresiva de ruido, la estimación de densidad penalizó las distorsiones estructurales introducidas por dicho procesamiento y favoreció los algoritmos de realce conservadores. Asimismo, se demostró que la métrica objetiva compuesta fracasó al predecir la calidad subjetiva del habla sintetizada, en tanto que la verosimilitud estimada por el modelo de difusión exhibió una correlación lineal y estadísticamente significativa. Se concluye que la elección del criterio de evaluación y del algoritmo de procesamiento debe subordinarse al objetivo del sistema final, y que los conjuntos de datos no profesionales resultan indispensables para capturar la variabilidad prosódica y la representatividad dialectal del español de Argentina.

\noindent\textbf{Palabras clave:} síntesis de voz, conjuntos de datos del habla, preprocesamiento de audio, lenguajes de bajos recursos
